% Compile in this order:
% pdflatex calabi-yau.tex
% asy calabi-yau-1.asy
% pdflatex calabi-yau.tex

\documentclass[border=2pt]{standalone}
\usepackage{asymptote}

\begin{document}
\begin{asy}[width=10cm,height=10cm]
settings.prc=false;
settings.render=8;
settings.twosided=true;
import graph3;

size3(200);
currentprojection=orthographic(3,3,2);

// Main light plus weaker fill lights, all strictly grayscale.
// Fill lights reveal curvature on faces pointing away from the main light.
currentlight=light(
  diffuse=new pen[] {
    gray(0.75), gray(0.35), gray(0.22), gray(0.18)
  },
  position=new triple[] {
    (8,10,2), (-6,3,5), (-4,-6,2), (2,-2,-6)
  }
);
currentlight.background=white;

// Soft satin finish: lift deep shadows while retaining local highlights.
material satin=material(
  diffusepen=gray(0.85),
  emissivepen=gray(0.21),
  specularpen=gray(0.65),
  shininess=0.50,
  fresnel0=0.08
);

int k1,k2,n=5;
real alpha=0.3*pi;

triple cy(pair z) {
  pair z1,z2;
  if(z==(0,0)) {
    z1=exp(2*pi*I*k1/n);
    z2=0;
  } else {
    z1=exp(2*pi*I*k1/n)*exp(log(cos(I*z))*2/n);
    z2=exp(2*pi*I*k2/n)*exp(log(-I*sin(I*z))*2/n);
  }
  return (
    z2.x,
    cos(alpha)*z1.y+sin(alpha)*z2.y,
    z1.x
  );
}

// Denser sampling reduces visible faceting without spline overshoot.
for(k1=0;k1<n;++k1) {
  for(k2=0;k2<n;++k2) {
    surface s=surface(cy,(-1,0),(1,0.5*pi),60,60);
    draw(s,surfacepen=satin,meshpen=nullpen);
  }
}
\end{asy}
\end{document}
